% VI.04.P4
\subsection*{Problem VI.04.P4 --- A bijective morphism that is not an isomorphism}
\label{prob:vi04-p4}
\paragraph{Problem.}
Let \(C=V(y^2-x^3)\) and \(\nu(t)=(t^2,t^3)\).  Prove that \(\nu:\mathbb A^1\to C\) is
bijective on \(k\)-rational points but is not an isomorphism of affine algebraic sets.
\ifdefined\FullProblemDossiers
\paragraph{Why this problem matters.}  It separates set-theoretic bijectivity from
algebraic isomorphism and explains why the warning ``bijective morphism need not be an
isomorphism'' is substantive.
\paragraph{Guided hints.}  If \((a,b)\in C\) and \(a\neq0\), try \(t=b/a\).  On rings,
compute the image of \(\nu^*\).
\paragraph{Solution.}
If \(a=0\), then \(b^2=0\), so \((a,b)=(0,0)=\nu(0)\).  If \(a\neq0\), put \(t=b/a\).
Then
\[
t^2=b^2/a^2=a,
\qquad
t^3=ta=b,
\]
so \(\nu\) is surjective; injectivity follows because \(t^2=s^2\) and \(t^3=s^3\) imply
\(t=s\) (including the zero case).  However
\[
\nu^*:k[x,y]/(y^2-x^3)\to k[t],\qquad x\mapsto t^2,\ y\mapsto t^3
\]
has image \(k[t^2,t^3]\), which does not contain \(t\).  Hence it is not surjective, so
\(\nu\) is not an isomorphism.
\paragraph{What happened mathematically.}  The missing function \(t\) detects the singular
point even though the underlying rational-point sets correspond bijectively.
\paragraph{Extensions.}  Compare this with \(t\mapsto(t,t^2)\) for the parabola.  Later,
normalization will make this cusp map conceptually canonical.
\paragraph{Provenance.}  New synthesis problem motivated by the chapter's isomorphism
criterion and the contrast with the legacy smooth-cubic problem.
\fi
